\chapter{Compiler Architecture}
\label{ch:compiler}

This chapter describes the architecture and implementation of the \haira{} compiler.

\section{Overview}

The \haira{} compiler is written in Go and transforms \code{.haira} source files into native binaries by generating Go code and invoking the Go toolchain. Agentic constructs --- \keyword{provider}, \keyword{tool}, \keyword{agent}, and \keyword{workflow} with \code{@} decorator triggers --- are first-class elements of the compilation pipeline.

The compiler pipeline:

\begin{lstlisting}[language={},numbers=none]
Source (.haira)
       |
       v
   +--------+
   | Lexer  |  --> Tokens
   +--------+
       |
       v
   +--------+
   | Parser |  --> AST
   +--------+
       |
       v
   +----------+
   | Resolver |  --> Merged AST (multi-file imports resolved)
   +----------+
       |
       v
   +----------+
   | Checker  |  --> Type info + diagnostics
   +----------+
       |
       v
   +----------+
   | Codegen  |  --> Go source code
   +----------+
       |
       v
   +----------+
   | go build |  --> Native binary
   +----------+
\end{lstlisting}

The key design decision: \haira{} compiles \emph{to Go}, not to machine code directly. This means every \haira{} program benefits from Go's battle-tested runtime (goroutines, garbage collector, HTTP stack, TLS) without reimplementing any of it. The generated Go code is readable and can be inspected via \code{haira emit}.

\begin{note}
There is no ``AI Phase'' in the pipeline. All LLM interaction happens at \emph{runtime} through agent method calls. The compiler validates, type-checks, and generates schemas for agentic constructs, then emits runtime library calls.
\end{note}

\section{Lexer}

The lexer is a hand-written scanner in \code{compiler/internal/lexer/} that converts UTF-8 source text into a stream of tokens.

\subsection{Token Types}

Tokens are defined in \code{compiler/internal/token/}:

\begin{lstlisting}[language={},numbers=none]
// Literals
Int, Float, String, InterpolatedString, TripleQuoteString

// Keywords (agentic)
Provider, Tool, Agent, Workflow

// Keywords (core)
Fn, Struct, Enum, Type, Import, Export, From, Pub,
If, Else, For, While, Match, Return, Break, Continue,
True, False, None, Some, And, Or, Not, In,
Spawn, Select, Try, Catch, Defer, Errdefer,
Step, Test, Assert, Let, Const

// Keywords (reserved for future)
Trait, Impl, Async

// Operators
Plus, Minus, Star, Slash, Percent,
Eq, EqEq, Ne, Lt, Gt, Le, Ge,
PlusEq, MinusEq, StarEq, SlashEq, PercentEq,
Pipe, PipeArrow, Amp, AmpEq, Caret, CaretEq, Tilde,
Shl, ShlEq, Shr, ShrEq, PipeEq,
Arrow, FatArrow, DotDot, DotDotEq, Dot, Question, At, Ellipsis

// Delimiters
LParen, RParen, LBracket, RBracket, LBrace, RBrace, Comma, Colon
\end{lstlisting}

\subsection{Lexer Features}

\begin{itemize}
    \item \textbf{Numeric literals} --- Decimal, hex (\code{0xFF}), binary (\code{0b1010}), octal (\code{0o77}), with underscores (\code{1\_000\_000})
    \item \textbf{String interpolation} --- \code{\$\{expr\}} syntax with nested brace tracking
    \item \textbf{Triple-quoted strings} --- \code{"""..."""} with automatic leading-whitespace dedenting
    \item \textbf{Nested block comments} --- \code{/* ... /* nested */ ... */}
    \item \textbf{Escape sequences} --- \code{\\n}, \code{\\t}, \code{\\r}, \code{\\\\}, \code{\\"}, \code{\\\{}, \code{\\\}}
\end{itemize}

\section{Parser}

The parser in \code{compiler/internal/parser/} uses recursive descent for statements and declarations, combined with Pratt (precedence-climbing) parsing for expressions.

\subsection{AST Nodes}

The AST is defined in \code{compiler/internal/ast/} using Go interfaces as sum types. All nodes carry a \code{Span} for error reporting.

\textbf{Top-level declarations:}

\begin{itemize}
    \item \code{TypeDef} --- struct with typed fields and optional defaults
    \item \code{EnumDef} --- enum with variants (optionally carrying fields)
    \item \code{TypeAlias} --- \code{type Name = Type}
    \item \code{FunctionDef} --- \code{fn name(params) -> Type \{ body \}}
    \item \code{MethodDef} --- \code{Type.method(params) -> Type \{ body \}} with implicit \code{self}
    \item \code{ImportDecl} --- four forms: basic, aliased, selective, glob
    \item \code{ExportDecl} --- \code{export \{ Name1, Name2 \}}
    \item \code{ProviderDecl} --- \code{provider name \{ key: value \}}
    \item \code{ToolDecl} --- tool with params, return type, triple-quoted description, body
    \item \code{AgentDecl} --- \code{agent name \{ key: value \}}
    \item \code{WorkflowDecl} --- workflow with optional trigger decorator, params, lifecycle hooks, body
    \item \code{TestDecl} --- \code{test "name" \{ body \}}
\end{itemize}

\textbf{Expressions (27 types):} Literals, identifiers, binary/unary ops, calls, method calls, field access, indexing, pipe, lambda, match, if-expressions, lists, maps, struct instances, ranges, error propagation (\code{?}), \code{orelse}, \code{some}/\code{none}, spawn, select, async blocks.

\textbf{Statements (14 types):} Assignment, let/const, if/else, for-in, while, return, try/catch, defer/errdefer, match, break/continue, step, assert, expression statements.

\textbf{Patterns (6 types):} Wildcard (\code{\_}), literal, identifier, constructor, or-pattern (\code{A | B}), range (\code{1..5}).

\subsection{Operator Precedence}

The Pratt parser implements 14 precedence levels (lowest to highest):

\begin{table}[h]
\centering
\begin{tabular}{cl}
\toprule
\textbf{Level} & \textbf{Operators} \\
\midrule
1 & \code{orelse} \\
2 & \code{|>} (pipe) \\
3 & \code{or} (logical) \\
4 & \code{and} (logical) \\
5 & \code{|} (bitwise OR) \\
6 & \code{\^{}} (bitwise XOR) \\
7 & \code{\&} (bitwise AND) \\
8 & \code{==}, \code{!=} \\
9 & \code{<}, \code{>}, \code{<=}, \code{>=}, \code{..}, \code{..=} \\
10 & \code{<<}, \code{>>} \\
11 & \code{+}, \code{-} \\
12 & \code{*}, \code{/}, \code{\%} \\
13 & \code{-}, \code{not}, \code{\~{}} (unary) \\
14 & \code{()}, \code{[]}, \code{.}, \code{?} (postfix) \\
\bottomrule
\end{tabular}
\caption{Operator precedence (lowest to highest)}
\end{table}

\section{Resolver}

The resolver in \code{compiler/internal/resolver/} handles multi-file module resolution:

\begin{enumerate}
    \item Reads the main source file and extracts import declarations.
    \item Resolves each import: standard library (built-in), project-local (\code{path/file.haira} or \code{path/mod.haira}), or external.
    \item Recursively resolves transitive imports.
    \item Detects circular dependencies (compile-time error).
    \item Filters exports: only \code{pub} items and agentic declarations are visible to importers.
    \item Produces a \code{Program} with a merged item list in dependency order.
\end{enumerate}

\begin{lstlisting}[language={},numbers=none]
// Resolution order for merged items:
// 1. Transitive module items (deepest dependencies first)
// 2. Direct import items (in source order)
// 3. Main file items
\end{lstlisting}

\section{Checker}

The type checker in \code{compiler/internal/checker/} performs semantic analysis in two passes:

\textbf{Pass 1 (Registration):} Walk all top-level items and register:
\begin{itemize}
    \item Struct types with field types
    \item Enum types with variants
    \item Function signatures
    \item Provider, tool, agent, workflow declarations
    \item Global variables (inferred from initializers)
\end{itemize}

\textbf{Pass 2 (Validation):} Check function/method bodies:
\begin{itemize}
    \item Type inference for expressions (literals, binary ops, calls, field access, indexing)
    \item Variable definition tracking (warn on undefined variables)
    \item Const immutability enforcement
    \item Agent field validation (model references a provider, tools exist, known fields)
    \item Enum exhaustiveness warnings in match expressions
    \item Return type checking
    \item Method \code{self} binding
\end{itemize}

\subsection{Type System}

\begin{lstlisting}[language={},numbers=none]
Primitive types: int, float, string, bool, any, void, error
Compound types: [T] (list), {K: V} (map)
Named types:    struct, enum
Function types: (T, U) -> V
\end{lstlisting}

\begin{note}
Generic types are parse-accepted but erased to \code{any} in v1. Full generics are planned for a future release.
\end{note}

\subsection{Agentic Validation}

The checker validates agentic declarations:

\begin{itemize}
    \item \textbf{Providers:} Warn on unknown fields. Valid fields: \code{model}, \code{api\_key}, \code{backend}, \code{host}, \code{temperature}, \code{max\_tokens}, \code{transport}, \code{command}, \code{args}, \code{url}.
    \item \textbf{Agents:} Require \code{model} field referencing a declared provider. Validate tool references. Check handoff targets exist. Warn on unknown fields.
    \item \textbf{Tools:} Enforce mandatory triple-quoted description (parser-level). Register parameters for schema generation.
    \item \textbf{Workflows:} Validate decorator triggers. Type-check body.
\end{itemize}

\section{Code Generation}

The code generator in \code{compiler/internal/codegen/} transforms the typed AST into Go source code. This is the core of \haira{}'s compilation strategy: generate idiomatic Go that calls into the \haira{} runtime library.

\subsection{Emission Strategy}

Code is emitted in 8 ordered passes to satisfy Go's declaration requirements:

\begin{enumerate}
    \item \textbf{Providers} --- \code{var Claude = \&haira.Provider\{...\}}
    \item \textbf{Tools} --- Tool functions + registration
    \item \textbf{Agents} --- Agent variables + initialization functions
    \item \textbf{Workflows} --- HTTP handler functions + SSE streaming handlers
    \item \textbf{Top-level variables} --- Global \code{var} declarations
    \item \textbf{Types and enums} --- Struct definitions, enum iota constants
    \item \textbf{Methods} --- Receiver methods (dot-attach syntax)
    \item \textbf{Functions} --- User functions (non-main)
    \item \textbf{Main function} --- Entry point with agent initialization in topological order
\end{enumerate}

\subsection{Agent Initialization Order}

Agents with handoff targets must be initialized after their targets. The codegen performs a topological sort on the handoff graph to determine initialization order:

\begin{lstlisting}
// Haira source:
agent Triage { handoffs: [Billing, Support] }
agent Billing { provider: Claude }
agent Support { provider: Claude }

// Generated init order:
// 1. Billing (no dependencies)
// 2. Support (no dependencies)
// 3. Triage  (depends on Billing, Support)
\end{lstlisting}

\subsection{Workflow Code Generation}

Workflows with triggers generate HTTP handlers. Streaming workflows generate both a \code{StreamHandler} (SSE) and a \code{Handler} (JSON fallback):

\begin{lstlisting}[language={},numbers=none]
// @webhook("/summarize")
// workflow Summarize(req: SumReq) -> Summary { ... }

// Generates:
// func SummarizeHandler(w, r)      -- JSON response
// func SummarizeStreamHandler(w, r) -- SSE streaming
// Route registration in main()
\end{lstlisting}

\subsection{Tool Schema Generation}

Each tool declaration generates a JSON schema from its parameter types:

\begin{table}[h]
\centering
\begin{tabular}{ll}
\toprule
\textbf{\haira{} Type} & \textbf{JSON Schema Type} \\
\midrule
\type{int} & \code{"integer"} \\
\type{float} & \code{"number"} \\
\type{bool} & \code{"boolean"} \\
\type{string} & \code{"string"} \\
\type{[T]} & \code{"array"} with \code{items} \\
\type{struct} & \code{"object"} with \code{properties} \\
\bottomrule
\end{tabular}
\caption{Type-to-JSON-schema mapping}
\end{table}

\subsection{Import Detection}

The codegen scans the AST to determine which Go imports are needed:

\begin{itemize}
    \item \code{fmt} --- string interpolation, print calls
    \item \code{encoding/json} --- JSON marshal/unmarshal calls
    \item \code{sync} --- spawn blocks (WaitGroup)
    \item \code{haira} --- runtime library (agents, providers, tools, workflows, HTTP)
\end{itemize}

\subsection{Go Compilation}

After generating Go source code, the codegen:

\begin{enumerate}
    \item Creates a temporary Go project directory with \code{go.mod}
    \item Extracts the embedded runtime bundle (\code{bundle.tar.gz}) containing the haira Go package
    \item Writes the generated \code{main.go}
    \item Runs \code{go build} to produce a native binary
    \item Moves the binary to the output path and cleans up
\end{enumerate}

\section{Runtime Library}
\label{sec:runtime}

The \haira{} runtime is a Go package (\code{package haira}) bundled into the compiler binary. It provides the execution environment for all agentic constructs.

\subsection{Runtime Structure}

The runtime is split into two layers:

\begin{itemize}
    \item \textbf{Primitive} (\code{primitive/haira/}) --- Core runtime: agent execution, provider management, tool registry, workflow orchestration, HTTP server, memory, I/O, strings, arrays, maps, math, JSON, regex, file system, time, environment variables.
    \item \textbf{Stdlib} (\code{stdlib/haira/}) --- External integrations: PostgreSQL, SQLite, vector databases, Excel, Slack, GitHub, GitLab.
\end{itemize}

Both layers use \code{package haira} and are merged at bundle time into a single Go package, embedded in the compiler binary as \code{bundle.tar.gz}.

\subsection{Provider Management}

The runtime manages LLM provider connections:

\begin{itemize}
    \item API key resolution from environment variables at startup
    \item Model routing based on provider configuration
    \item HTTP connection pooling to provider APIs
    \item Support for Anthropic, OpenAI, Ollama, and MCP (Model Context Protocol) transports
\end{itemize}

\subsection{Agent Execution Loop}

\begin{enumerate}
    \item Assemble message payload (system prompt, conversation history, user message).
    \item Send request to the LLM provider.
    \item If response contains tool calls, execute each tool and collect results.
    \item Send tool results back to the LLM.
    \item Repeat steps 2--4 until final response or \code{max\_steps} reached.
    \item For \code{.run()} calls, parse JSON response into the annotated output type.
    \item For \code{.stream()} calls, yield tokens incrementally via SSE.
\end{enumerate}

\subsection{Workflow Orchestration}

\begin{itemize}
    \item \textbf{HTTP server} --- Built on Go's \code{net/http}. Handles \code{@webhook} (POST/GET/PUT/DELETE) and \code{@websocket} triggers.
    \item \textbf{Cron scheduler} --- Executes \code{@cron}-triggered workflows on 5-field UNIX cron schedules.
    \item \textbf{Event bus} --- In-process event dispatch for \code{@event}-triggered workflows via \code{haira.emit()}.
    \item \textbf{Step telemetry} --- Named \code{step} blocks emit timing and status events for observability.
    \item \textbf{Run tracking} --- Each workflow execution gets a unique run ID with step-level SSE streaming.
\end{itemize}

\subsection{Session Management}

\begin{itemize}
    \item Each session maintains an isolated conversation history.
    \item \code{conversation(max\_turns: N)} memory trims to the last N message pairs.
    \item Sessions are stored in-memory by default, with optional persistence via the Store interface.
\end{itemize}

\section{Project Structure}

The compiler is organized as a Go module:

\begin{lstlisting}[language={},numbers=none]
compiler/
  main.go                          # CLI entry point
  internal/
    token/                         # Token type definitions
    lexer/                         # Hand-written scanner
    ast/                           # AST node definitions
    parser/                        # Recursive descent + Pratt parser
    resolver/                      # Multi-file import resolution
    checker/                       # Type checking + semantic analysis
    codegen/                       # Go code generation
      codegen.go                   #   8-pass emission orchestrator
      emitter.go                   #   Output buffer + indentation
      expressions.go               #   Expression-to-Go conversion
      statements.go                #   Statement emission
      agentic.go                   #   Agent/tool/provider/workflow emission
      stdlib.go                    #   Standard library function mapping
      types.go                     #   Haira-type-to-Go-type conversion
      project.go                   #   Temp project setup + go build
      test_codegen.go              #   Test harness generation
    driver/                        # Pipeline orchestration
    errors/                        # Diagnostic system (pretty-printing)
    formatter/                     # Source code formatter
    lsp/                           # Language server protocol
    manifest/                      # package.haira handling
    runtime/                       # Embedded runtime bundle (bundle.tar.gz)

primitive/haira/                   # Core runtime (Go)
stdlib/haira/                      # External integrations (Go)
ui/sdk/                            # Web component SDK (TypeScript)
ui/application/                    # Default HTML templates
\end{lstlisting}

\begin{table}[h]
\centering
\begin{tabular}{ll}
\toprule
\textbf{Package} & \textbf{Responsibility} \\
\midrule
\code{token} & Token kind constants and Token struct \\
\code{lexer} & Tokenizes source text \\
\code{ast} & AST node types (Go interfaces as sum types) \\
\code{parser} & Builds AST from tokens \\
\code{resolver} & Resolves imports and produces merged program \\
\code{checker} & Type checking, inference, and semantic validation \\
\code{codegen} & Generates Go source code from typed AST \\
\code{driver} & Orchestrates the full compilation pipeline \\
\code{errors} & Diagnostic formatting with source context \\
\code{formatter} & Haira source code formatting \\
\code{lsp} & Language server (hover, completion, go-to-definition) \\
\code{manifest} & Package manifest (package.haira) support \\
\bottomrule
\end{tabular}
\caption{Compiler package responsibilities}
\end{table}

\section{CLI Commands}

\begin{lstlisting}[language=bash,numbers=none]
haira build [file] [-o output]   # Compile to native binary
haira run [file]                  # Compile and execute
haira check [file]                # Type-check without codegen
haira emit [file]                 # Show generated Go code
haira test [file]                 # Run test blocks
haira fmt [file]                  # Format source in-place
haira parse [file]                # Show the AST
haira lex [file]                  # Show tokens
haira init                        # Create package.haira manifest
haira lsp                         # Start language server (stdio)
haira version                     # Show version
\end{lstlisting}

If no file is specified, \code{haira} looks for \code{package.haira} in the current directory and uses its entry point.

\section{Error Messages}

The compiler produces contextual error messages with source location, underlines, and help text.

\subsection{Missing Tool Description}

\begin{lstlisting}[language={},numbers=none]
error: tool is missing a description
  --> tools.haira:5:1
   |
5  | tool search(query: string) -> [Result] {
   | ^^^^^^^^^^^^^^^^^^^^^^^^^^^^^^^^^^^^^^^^
   |
   = help: add a triple-quoted description: """Search the web"""
   = note: tool descriptions are sent to the LLM and cannot be omitted
\end{lstlisting}

\subsection{Agent Referencing Undefined Provider}

\begin{lstlisting}[language={},numbers=none]
error: undefined provider 'openai'
  --> agents.haira:8:12
   |
8  |     provider: openai
   |               ^^^^^^ not found in this scope
   |
   = help: declare a provider with: provider openai { model: "gpt-4o" ... }
\end{lstlisting}

\subsection{Undefined Variable}

\begin{lstlisting}[language={},numbers=none]
warning: undefined variable 'foo'
  --> main.haira:10:5
   |
10 |     io.println(foo)
   |                ^^^
   |
   = hint: did you mean to declare it?
\end{lstlisting}

\subsection{Workflow with Invalid Trigger}

\begin{lstlisting}[language={},numbers=none]
error: unknown trigger '@daily'
  --> workflows.haira:12:1
   |
12 | @daily
   | ^^^^^^ not a recognized trigger
   |
   = help: valid triggers: @webhook, @websocket, @cron, @event, @manual
   = help: for daily execution, use: @cron("0 0 * * *")
\end{lstlisting}

\section{Runtime Bundling}

The runtime is bundled into the compiler binary at build time:

\begin{enumerate}
    \item \code{make bundle-runtime} copies \code{primitive/haira/*.go} and \code{stdlib/haira/*.go} into a staging directory.
    \item The UI SDK is built (\code{bun build}) and its output included.
    \item Everything is archived as \code{bundle.tar.gz}.
    \item The archive is embedded in the compiler binary via Go's \code{go:embed} directive.
    \item At compile time, the codegen extracts this bundle into a temporary Go project, writes the generated \code{main.go}, and runs \code{go build}.
\end{enumerate}

This approach produces fully standalone binaries: the compiled \haira{} program contains everything it needs, including the full runtime library. No separate runtime installation is required.
